\subsection{Gorenstein rings}

DEFINITION 1.--- A ring A is said to be a Gorenstein ring if it is Noetherian and the \(\mathrm{A}_{\mathfrak{m}}\) -module \(\mathrm{A}_{\mathfrak{m}}\) has finite injective dimension for every maximal ideal \(\mathfrak{m}\) of A.

For a local Noetherian ring A to be a Gorenstein ring, it is necessary and sufficient that \(\mathrm{di}_{\mathrm{A}}(\mathrm{A})\) be finite; for a Noetherian ring A to be a Gorenstein ring, it is necessary and sufficient that this be so for \(\Lambda_{m}\) for every maximal ideal m of A.

PROPOSITION 10.- Let A be a Gorenstein ring; then A is a Macaulay ring, and satisfies \(\mathrm{di}_{\Lambda}(\mathrm{A}) = \dim(\mathrm{A})\) .

For every maximal ideal m of A, one has

\[
\begin{array}{l l} \dim (A _ {m}) & \leqslant \mathrm{di} _ {A _ {m}} (A _ {m}) \\ \mathrm{di} _ {A _ {m}} (A _ {m}) & = \operatorname{prof} (A _ {m}) \\ \operatorname{prof} (A _ {m}) & \leqslant \dim (A _ {m}) \end{array} \quad \begin{array}{l l} & (n ^ {\circ} 6, \text { prop. } 8) \\ & (n ^ {\circ} 6, \text { prop. } 9) \\ & (\S 1, n ^ {\circ} 4, \text { cor. } 2 \text { du   th. } 2); \end{array}
\]

whence it follows that A is a Macaulay ring, and that \(\mathrm{di}_{\mathrm{A}}(\mathrm{A}) = \dim(\mathrm{A})\) by passing to the supremum (no. 2, prop. 3).

Thus the Noetherian rings A such that \(\mathrm{di}_{\mathrm{A}}(\mathrm{A})\) is finite are the Gorenstein rings of finite dimension (prop. 3 of no. 2), and the Noetherian rings such that the A-module A is injective are the Artinian Gorenstein rings.

Examples. 1) For every multiplicative set S of a Gorenstein ring A, the ring of fractions \(S^{-1}A\) is a Gorenstein ring: indeed, let q be a maximal ideal of \(S^{-1}A\) ; it is of the form \(S^{-1}p\) , where p is a prime ideal of A not meeting S. Let m be a maximal ideal of A containing p; then the ring \(B = (S^{-1}A)_q\) is isomorphic to \(A_p\) (II, § 2, no. 5, prop. 11), hence to a ring of fractions of \(A_m\) , and consequently satisfies \(\text{di}_B(B) < +\infty\) (no. 2, cor. 1 of prop. 3).

\begin{enumerate}
\def\labelenumi{\arabic{enumi})}
\setcounter{enumi}{1}
\item
  Let A be a Gorenstein ring and let J be an ideal of A, generated by an A-regular sequence x. The quotient ring A/J is a Gorenstein ring: indeed, for every maximal ideal m of A containing J, the image in \(A_{m}\) of the sequence x is \(A_{m}\) -regular and generates the ideal \(J_{m}\) , so that \(A_{m}/J_{m}\) is a Gorenstein ring by the cor. of prop. 7 (no. 4).
\item
  Let A be a local Noetherian ring, J an ideal of A generated by an A-regular sequence of elements of \(\mathfrak{m}_{\mathrm{A}}\) . If A/J is a Gorenstein ring, then so is A (no. 4, cor. of prop. 7).
\item
  Let A be a regular local Noetherian ring; then A is a Gorenstein ring. Indeed, let x be a system of coordinates of A (VIII, § 5, no. 1, def. 1). The sequence x is A-regular (loc. cit., no. 2, th. 1) and generates the ideal m\_A; one can therefore apply example 3.
\item
  Every quotient ring of a principal ideal ring is a Gorenstein ring (example 2). In particular, every monogenic algebra over a field is a Gorenstein ring.
\end{enumerate}

Lemma 1.- Let A be a local Artinian ring. The following conditions are equivalent:

\begin{enumerate}
\def\labelenumi{(\roman{enumi})}
\item
  A is a Gorenstein ring;
\item
  the A-module A is injective;
\item
  the \(\kappa_{\mathrm{A}}\) -vector space \(\operatorname{Hom}_{\mathrm{A}}(\kappa_{\mathrm{A}}, \mathrm{A})\) has dimension 1.
\end{enumerate}

Recall (A, VIII, p.~3.5) that A possesses a nonzero minimal ideal; such an ideal is a simple module, hence isomorphic to \(\kappa_{\mathrm{A}}\) . The \(\kappa_{\mathrm{A}}\) -vector space \(\mathrm{Hom}_{\mathrm{A}}(\kappa_{\mathrm{A}}, \Lambda)\) is therefore nonzero; to say that it has dimension 1 means that A contains only one nonzero minimal ideal, which is then the socle of A (A, VIII, § 4, no. 6).

The equivalence of (i) and (ii) was proved after prop. 10. Suppose the A-module A is injective. Let \(x\) and \(y\) be two nonzero elements of A annihilated by \(\mathfrak{m}_{\Lambda}\) . There exists a unique A-linear map \(\varphi : Ax \to A\) such that \(\varphi(x) = y\) ; since A is injective, it extends to an endomorphism of \(\Lambda\) , which implies that \(y\) belongs to \(Ax\) , whence (iii).

Suppose conversely that \(\mathrm{Hom}_{\mathrm{A}}(\kappa_{\mathrm{A}},\Lambda)\) has dimension 1. Let M be a finitely generated A-module; it has finite length, hence possesses a composition series whose quotients are isomorphic to \(\kappa_{\mathrm{A}}\) . From this one deduces by induction on the length of M the inequality \(\mathrm{long}_{\mathrm{A}}(\mathrm{Hom}_{\mathrm{A}}(\mathrm{M},\Lambda)) \leqslant \mathrm{long}_{\mathrm{A}}(\mathrm{M})\) . In the exact sequence of extension modules

\[
0 \to \operatorname{Hom} _ {\mathrm{A}} (\kappa_ {\mathrm{A}}, \mathrm{A}) \longrightarrow \operatorname{Hom} _ {\mathrm{A}} (\mathrm{A}, \mathrm{A}) \stackrel {{\alpha}} {{\longrightarrow}} \operatorname{Hom} _ {\mathrm{A}} (\mathfrak {m} _ {\mathrm{A}}, \mathrm{A}) \longrightarrow \operatorname{Ext} _ {\mathrm{A}} ^ {1} (\kappa_ {\Lambda}, \mathrm{A}) \to 0,
\]

one therefore has \(\operatorname{long}_{\mathrm{A}}(\operatorname{Hom}_{\mathrm{A}}(\mathfrak{m}_{\mathrm{A}},\mathrm{A})) \leqslant \operatorname{long}_{\mathrm{A}}(\mathfrak{m}_{\mathrm{A}}) = \operatorname{long}(\mathrm{A}) - 1 = \operatorname{long}_{\mathrm{A}}(\operatorname{Im}\alpha)\) . Consequently \(\alpha\) is surjective, \(\operatorname{Ext}_{\mathrm{A}}^{1}(\kappa_{\mathrm{A}},\mathrm{A})\) is zero, and the A-module A is injective (no. 3, prop. 4).

Lemma 2.--- Let \(\Lambda\) be a local Noetherian ring such that \(\mathrm{di}_{\mathbf{A}}(\mathbf{A}) = +\infty\) . One has \(\operatorname{Ext}_{\mathbf{A}}^{i}(\kappa_{\Lambda}, \mathbf{A}) \neq 0\) for every integer \(i \geqslant \dim(\mathbf{A})\) .

Reason by induction on \(\dim(A)\) . When \(\dim(A) = 0\) , \(\mathfrak{m}_A\) is the unique prime ideal of \(A\) and the assertion follows from prop. 6 of no. 3. Suppose therefore \(\dim(A) > 0\) and let \(\mathfrak{p}\) a prime ideal of \(A\) distinct from \(\mathfrak{m}_A\) ; we then have \(\dim(A_{\mathfrak{p}}) < \dim(A)\) If \(\mathrm{di}_{A_{\mathfrak{p}}}(A_{\mathfrak{p}}) = +\infty\) , the induction hypothesis implies \(\operatorname{Ext}_{A_{\mathfrak{p}}}^{j}(\kappa(\mathfrak{p}), A_{\mathfrak{p}}) \neq 0\) for every integer \(j \geqslant \dim(A_{\mathfrak{p}})\) ; by lemma 3 of § 1, no. 7, this implies \(\operatorname{Ext}_{A}^{i}(\kappa_{A}, A) \neq 0\) for every integer \(i \geqslant \dim(A_{\mathfrak{p}}) + \dim(A/\mathfrak{p})\) and in particular for \(i \geqslant \dim(A)\) (VIII, § 1, no. 3, prop. 8, b)).

We still have to deal with the case where \(\mathrm{di}_{\mathrm{A}}(\mathrm{A})\) is infinite but where the injective dimension of \(A_{p}\) is finite for every prime ideal p of A distinct from \(m_{A}\) . For such an ideal, we have in this case, by prop. 10, \(\mathrm{di}_{\mathrm{A}_{\mathfrak{p}}}(\mathrm{A}_{\mathfrak{p}})=\dim(\mathrm{A}_{\mathfrak{p}})<\dim(\mathrm{A})\) , hence \(\operatorname{Ext}_{\mathrm{A}_{\mathfrak{p}}}^{i}(\kappa(\mathfrak{p}),\mathrm{A}_{\mathfrak{p}})=0\) for \(i\geqslant\dim(\Lambda)\) . Since \(\mathrm{di}_{\mathrm{A}}(\mathrm{A})\) is infinite, prop. 6 of no. 3 then imposes \(\operatorname{Ext}_{\mathrm{A}}^{i}(\kappa_{\mathrm{A}},\mathrm{A})\neq0\) for every integer \(i\geqslant\dim(\mathrm{A})\) .

THEOREM 2 (Bass). Let A be a Noetherian local ring; set \(d = \dim(A)\) . Let \(\mathbf{x} = \cdot(x_1, \ldots, x_d)\) a maximal secant sequence of elements of \(\mathfrak{m}_A\) , and by \(x\) the ideal it generates. The following conditions are equivalent:

\begin{enumerate}
\def\labelenumi{(\roman{enumi})}
\item
  A is a Gorenstein ring;
\item
  one has \(\mathrm{di}_{\mathrm{A}}(\mathrm{A}) = d\)
\item
  there exists an integer \(i>d\) such that \(\operatorname{Ext}_{\mathrm{A}}^{i}(\kappa_{\mathrm{A}},\mathrm{A})=0\) ;
\item
  \(\operatorname{Ext}_{\mathrm{A}}^{i}(\kappa_{\mathrm{A}},\mathrm{A}) = 0\) for \(i < d\) and the \(\kappa_{\mathrm{A}}\) -vector space \(\operatorname{Ext}_{\mathrm{A}}^{d}(\kappa_{\mathrm{A}},\mathrm{A})\) is of dimension 1;
\item
  the ring A is Macaulay and the \(\kappa_{\mathrm{A}}\) -vector space \(\operatorname{Hom}_{\Lambda}(\kappa_{\mathrm{A}}, \mathrm{A}/x)\) is of dimension 1 ;
\item
  the sequence \(\mathbf{x}\) is A-regular and the \(\kappa_{\mathrm{A}}\) -vector space \(\operatorname{Hom}_{\mathrm{A}}(\kappa_{\mathrm{A}},\mathrm{A} / x)\) has dimension 1.
\end{enumerate}

The equivalence of (i), (ii), and (iii) follows from Lemma 2 and Prop. 10. If \(\mathrm{Ext}_{\mathrm{A}}^{i}(\kappa_{\mathrm{A}},\mathrm{A})\) is zero for every integer \(i < d\) , the ring A is Macaulay (§ 2, No. 3, prop. 3); if the ring A is Macaulay, the sequence x is A-regular (loc. cit., th. 1); if the sequence x is A-regular, we have \(\mathrm{Ext}_{\mathrm{A}}^{i}(\kappa_{\mathrm{A}},\mathrm{A}) = 0\) for \(i < d\) and the \(\kappa_{\mathrm{A}}\) -vector spaces \(\mathrm{Ext}_{\mathrm{A}}^{d}(\kappa_{\mathrm{A}},\mathrm{A})\) and \(\mathrm{Hom}_{\mathrm{A}}(\kappa_{\mathrm{A}},\mathrm{A}/x)\) are isomorphic (A, X, p.~166, prop. 9). This proves the equivalence of conditions (iv), (v), and (vi).

Let us prove that (i) implies (v): if A is a Gorenstein ring, it is a Macaulay ring (prop. 10). The sequence x is then A-regular (§ 2, no. 3, prop. 4), hence A/x is an Artinian Gorenstein ring (example 2), so that the \(\kappa_{\mathrm{A}}\) -vector space \(\mathrm{Hom}_{\mathrm{A}}(\kappa_{\mathrm{A}}, \mathrm{A}/x)\) has dimension 1.

Finally, let us prove that (vi) implies (i): under hypothesis (vi), the ring A/x is a Gorenstein ring (lemma 1), hence A is a Gorenstein ring (example 3).

COROLLARY.--- Let Λ be a Noetherian local ring of dimension d. The Λ-module \(\operatorname{Ext}_{\mathrm{A}}^{d}(\mathbf{k}_{\mathrm{A}},\Lambda)\) is not zero.

This follows from Thm. 2 if \(\Lambda\) is a Gorenstein ring and from Lemma 2 otherwise.

PROPOSITION 11.--- Let A be a Noetherian ring. The following conditions are equivalent:

\begin{enumerate}
\def\labelenumi{(\roman{enumi})}
\item
  A is a Gorenstein ring;
\item
  for every prime ideal \(\mathfrak{p}\) of A, the \(\kappa(\mathfrak{p})\) -vector space \(\operatorname{Ext}_{\Lambda_{\mathfrak{p}}}^{i}(\kappa(\mathfrak{p}),\mathrm{A}_{\mathfrak{p}})\) is zero for \(i \neq \operatorname{ht}(\mathfrak{p})\) and of dimension 1 for \(i = \operatorname{ht}(\mathfrak{p})\) .
\item
  for every maximal ideal \(\mathfrak{m}\) from A, there exists an integer \(i > \mathrm{ht}(\mathfrak{m})\) such that \(\operatorname{Ext}_{\Lambda_{\mathfrak{m}}}^{i}(\mathrm{A}/\mathfrak{m},\mathrm{A}_{\mathfrak{m}}) = 0\) .
\end{enumerate}

(i)⇒(ii): This follows from Theorem 2 applied to the Gorenstein local ring A. \(_{p}\) (example 1).

(ii)⇒(iii): this is trivial.

(iii)⇒(i): under hypothesis (iii), A \(_{m}\) is a Gorenstein ring for every maximal ideal m of A (th. 2), and A is Gorenstein.

